% QUESTAO
Sabemos que ao multiplicarmos uma matriz $A$, de ordem $n$, pela sua inversa $A^{-1}$, obtemos a matriz identidade $I_n$ (elemento neutro da multiplicação matricial), ou seja
\[A\cdot A^{-1}=I_n\textnormal{\quad e \quad}A^{-1}\cdot A=I_n \]
Partindo dessa definição, calcule a matriz $A^{-1}=\left[ \begin{array}{cc}
a   &  c \\
b   &  d
\end{array} \right]$, em que $A = \left[ \begin{array}{cc}
5  & -2\\
-2  &  1
\end{array} \right]$

% ALTERNATIVAS
$A^{-1} = \left[ \begin{array}{cc} 1   &  2\\ 2  &   5 \end{array} \right]$
$A^{-1} = \left[ \begin{array}{cc} -3  &  2\\ 2  & -1 \end{array} \right]$
$A^{-1} = \left[ \begin{array}{cc} 5  & -2\\ -2  &  1 \end{array} \right]$
$A^{-1} = \left[ \begin{array}{cc} 1   &  2\\ 1  &   1 \end{array} \right]$
$A^{-1} = \left[ \begin{array}{cc} 1  &  2\\ 2  &  3 \end{array} \right]$


